% GENERATED FILE. Edit canonical lessons, then run npm run generate:books.
% canonical-source-hash: fnv1a64:9801199edd723f6e
% canonical-lessons: TE-C456-sikaru, TE-C456-viharincu, TE-C456-drsyam, TE-C456-sandhya, TE-C456-majili

\chapter{Strolling Around}
\label{ch:te-strolling-around}
\hlchaptermodality{456}

\begin{chapteropening}
\textbf{By the end of this chapter:} \emph{I can say a jaunt, to stroll, a view, dusk and a stopover.}

Complete the last lesson of chapter 456: I can say a jaunt, to stroll, a view, dusk and a stopover.
\end{chapteropening}

\section[\texorpdfstring{\te{షికారు} (ṣikāru)}{షికారు (ṣikāru)}]{\te{షికారు} (ṣikāru) — a jaunt; an outing}

\label{lesson:TE-C456-sikaru}



\begin{quote}
Before the new one: say the Telugu for entry, then the Telugu for a pleasure trip.
\end{quote}

\subsection*{You'll want to know: \te{షికారు}}
\textbf{\te{షికారు}} — \emph{ṣikāru} — \textquotedblleft{}a jaunt; an outing\textquotedblright{}.

\begin{quote}
\te{సాయంత్రం} \te{మేము} \te{సముద్రం} \te{దగ్గర} \te{షికారుకు} \te{వెళ్ళాం}.
\end{quote}

\begin{quote}
\emph{sāyantraṁ mēmu samudraṁ daggara ṣikāruku veḷḷāṁ.}
\end{quote}

\begin{quote}
In the evening we went for a stroll by the sea.
\end{quote}

\te{షికారుకు} \te{వెళ్ళు}, to go for an outing. An outing with no aim but pleasure.

\begin{cousinweb}
\textbf{\te{షికారు}} (\emph{ṣikāru}) is a loan from Persian \emph{shikār}, \textquotedblleft{}hunting, the chase\textquotedblright{}, which reached Telugu through Hindustani; a hunting trip became any trip for pleasure.
\end{cousinweb}

\begin{grammarlens}[title={What you've built}]
One more word for getting around.
\end{grammarlens}

\subsection*{Your turn}
\begin{itemize}
\raggedright
  \item \emph{Say it:} \emph{ṣikāru}
  \item \emph{Say it:} \emph{ṣikāru}, once more
  \item \emph{Say it:} \emph{sāyantraṁ mēmu samudraṁ daggara ṣikāruku veḷḷāṁ.} — In the evening we went for a stroll by the sea.
  \item \emph{Recall:} say the Telugu for entry, then the Telugu for a pleasure trip, then say \emph{ṣikāru} again
\end{itemize}

\begin{tcolorbox}[breakable,colback=teal!4,colframe=teal!35!black,title={Before you move on}]
What does \te{షికారు} mean? (\textquotedblleft{}A jaunt; an outing\textquotedblright{}.) Say it once more.
\end{tcolorbox}
\section[\texorpdfstring{\te{విహరించు} (vihariñcu)}{విహరించు (vihariñcu)}]{\te{విహరించు} (vihariñcu) — to stroll; to roam for pleasure}

\label{lesson:TE-C456-viharincu}



\begin{quote}
Before the new one: say the Telugu for a pleasure trip, then the Telugu for a jaunt.
\end{quote}

\subsection*{You'll want to know: \te{విహరించు}}
\textbf{\te{విహరించు}} — \emph{vihariñcu} — \textquotedblleft{}to stroll; to roam for pleasure\textquotedblright{}.

\begin{quote}
\te{పర్యాటకులు} \te{తోటలో} \te{హాయిగా} \te{విహరించారు}.
\end{quote}

\begin{quote}
\emph{paryāṭakulu tōṭalō hāyigā vihariñcāru.}
\end{quote}

\begin{quote}
The tourists strolled happily in the garden.
\end{quote}

Past (I …): \te{విహరించాను}, \emph{vihariñcānu}.

A rather literary verb; \te{విహారయాత్ర}, an excursion, holds the same root.

\begin{cousinweb}
\textbf{\te{విహరించు}} (\emph{vihariñcu}) is built from Sanskrit \emph{vihāra}, \textquotedblleft{}roaming, pleasure\textquotedblright{} (CDIAL 12033), from \emph{vi-hṛ}, \textquotedblleft{}to roam about\textquotedblright{}, with the Telugu verb ending -\te{ఇంచు}.
\end{cousinweb}

\begin{grammarlens}[title={What you've built}]
One more word for getting around.
\end{grammarlens}

\subsection*{Your turn}
\begin{itemize}
\raggedright
  \item \emph{Say it:} \emph{vihariñcu}
  \item \emph{Say it:} \emph{vihariñcu}, once more
  \item \emph{Say it:} \emph{paryāṭakulu tōṭalō hāyigā vihariñcāru.} — The tourists strolled happily in the garden.
  \item \emph{Recall:} say the Telugu for a pleasure trip, then the Telugu for a jaunt, then say \emph{vihariñcu} again
\end{itemize}

\begin{tcolorbox}[breakable,colback=teal!4,colframe=teal!35!black,title={Before you move on}]
What does \te{విహరించు} mean? (\textquotedblleft{}To stroll; to roam for pleasure\textquotedblright{}.) Say it once more.
\end{tcolorbox}
\section[\texorpdfstring{\te{దృశ్యం} (dṛśyaṁ)}{దృశ్యం (dṛśyaṁ)}]{\te{దృశ్యం} (dṛśyaṁ) — a view; a scene}

\label{lesson:TE-C456-drsyam}



\begin{quote}
Before the new one: say the Telugu for a jaunt, then the Telugu for to stroll.
\end{quote}

\subsection*{You'll want to know: \te{దృశ్యం}}
\textbf{\te{దృశ్యం}} — \emph{dṛśyaṁ} — \textquotedblleft{}a view; a scene\textquotedblright{}.

\begin{quote}
\te{కొండ} \te{మీద} \te{నుంచి} \te{దృశ్యం} \te{చాలా} \te{అందంగా} \te{ఉంది}.
\end{quote}

\begin{quote}
\emph{koṇḍa mīda nuñci dṛśyaṁ cālā andaṅgā undi.}
\end{quote}

\begin{quote}
The view from the top of the hill is very beautiful.
\end{quote}

A view of nature, and also a scene in a film or a play.

\begin{cousinweb}
\textbf{\te{దృశ్యం}} (\emph{dṛśyaṁ}) is borrowed from Sanskrit \emph{dṛśya}, \textquotedblleft{}visible, what is to be seen\textquotedblright{} (Monier-Williams), from \emph{dṛś}, \textquotedblleft{}to see\textquotedblright{}.
\end{cousinweb}

\begin{grammarlens}[title={What you've built}]
One more word for getting around.
\end{grammarlens}

\subsection*{Your turn}
\begin{itemize}
\raggedright
  \item \emph{Say it:} \emph{dṛśyaṁ}
  \item \emph{Say it:} \emph{dṛśyaṁ}, once more
  \item \emph{Say it:} \emph{koṇḍa mīda nuñci dṛśyaṁ cālā andaṅgā undi.} — The view from the top of the hill is very beautiful.
  \item \emph{Recall:} say the Telugu for a jaunt, then the Telugu for to stroll, then say \emph{dṛśyaṁ} again
\end{itemize}

\begin{tcolorbox}[breakable,colback=teal!4,colframe=teal!35!black,title={Before you move on}]
What does \te{దృశ్యం} mean? (\textquotedblleft{}A view; a scene\textquotedblright{}.) Say it once more.
\end{tcolorbox}
\section[\texorpdfstring{\te{సంధ్య} (sandhya)}{సంధ్య (sandhya)}]{\te{సంధ్య} (sandhya) — dusk; twilight}

\label{lesson:TE-C456-sandhya}



\begin{quote}
Before the new one: say the Telugu for to stroll, then the Telugu for a view.
\end{quote}

\subsection*{You'll want to know: \te{సంధ్య}}
\textbf{\te{సంధ్య}} — \emph{sandhya} — \textquotedblleft{}dusk; twilight\textquotedblright{}.

\begin{quote}
\te{సంధ్య} \te{వేళ} \te{ఆకాశం} \te{ఎర్రగా} \te{ఉంది}.
\end{quote}

\begin{quote}
\emph{sandhya vēḷa ākāśaṁ erragā undi.}
\end{quote}

\begin{quote}
At dusk the sky is red.
\end{quote}

\te{సంధ్య} \te{వేళ}, at dusk. The word also names the meeting of day and night at dawn.

\begin{cousinweb}
\textbf{\te{సంధ్య}} (\emph{sandhya}) is borrowed from Sanskrit \emph{sandhyā}, \textquotedblleft{}a joining, the twilight\textquotedblright{} where day and night meet (Monier-Williams).
\end{cousinweb}

\begin{grammarlens}[title={What you've built}]
One more word for getting around.
\end{grammarlens}

\subsection*{Your turn}
\begin{itemize}
\raggedright
  \item \emph{Say it:} \emph{sandhya}
  \item \emph{Say it:} \emph{sandhya}, once more
  \item \emph{Say it:} \emph{sandhya vēḷa ākāśaṁ erragā undi.} — At dusk the sky is red.
  \item \emph{Recall:} say the Telugu for to stroll, then the Telugu for a view, then say \emph{sandhya} again
\end{itemize}

\begin{tcolorbox}[breakable,colback=teal!4,colframe=teal!35!black,title={Before you move on}]
What does \te{సంధ్య} mean? (\textquotedblleft{}Dusk; twilight\textquotedblright{}.) Say it once more.
\end{tcolorbox}
\section[\texorpdfstring{\te{మజిలీ} (majilī)}{మజిలీ (majilī)}]{\te{మజిలీ} (majilī) — a stopover; a halt}

\label{lesson:TE-C456-majili}



\begin{quote}
Before the new one: say the Telugu for a view, then the Telugu for dusk.
\end{quote}

\subsection*{You'll want to know: \te{మజిలీ}}
\textbf{\te{మజిలీ}} — \emph{majilī} — \textquotedblleft{}a stopover; a halt\textquotedblright{}.

\begin{quote}
\te{రెండు} \te{రోజుల} \te{ప్రయాణంలో} \te{మొదటి} \te{మజిలీ} \te{ఒక} \te{చిన్న} \te{పట్టణం}.
\end{quote}

\begin{quote}
\emph{reṇḍu rōjula prayāṇaṁlō modaṭi majilī oka cinna paṭṭaṇaṁ.}
\end{quote}

\begin{quote}
On the two-day journey the first stopover was a small town.
\end{quote}

A stop on a long journey. \te{మజిలీ} \te{కథలు} are stories told at each halt of a journey.

\begin{cousinweb}
\textbf{\te{మజిలీ}} (\emph{majilī}) is usually traced to Arabic \emph{manzil}, \textquotedblleft{}a place where travellers halt\textquotedblright{}, which came into Telugu through Persian and Hindustani.
\end{cousinweb}

\begin{grammarlens}[title={What you've built}]
One more word for getting around.
\end{grammarlens}

\subsection*{Your turn}
\begin{itemize}
\raggedright
  \item \emph{Say it:} \emph{majilī}
  \item \emph{Say it:} \emph{majilī}, once more
  \item \emph{Say it:} \emph{reṇḍu rōjula prayāṇaṁlō modaṭi majilī oka cinna paṭṭaṇaṁ.} — On the two-day journey the first stopover was a small town.
  \item \emph{Recall:} say the Telugu for a view, then the Telugu for dusk, then say \emph{majilī} again
\end{itemize}

\begin{tcolorbox}[breakable,colback=teal!4,colframe=teal!35!black,title={Before you move on}]
What does \te{మజిలీ} mean? (\textquotedblleft{}A stopover; a halt\textquotedblright{}.) Say it once more.
\end{tcolorbox}
